\chapter{Formato dei Dati di Input e Output}

\section{Input}

Il programma legge i dati da un file di testo passato come argomento 
da riga di comando:

\begin{lstlisting}
./percorsi input.txt
\end{lstlisting}

Il file di input ha il seguente formato:

\begin{itemize}
    \item \textbf{Prima riga}: due interi separati da uno spazio, 
    $N$ e $P$, dove $N$ è il numero di città (numerate da $1$ a $N$, 
    con $0$ che rappresenta la capitale) e $P$ è il numero di tratte;
    \item \textbf{Righe successive}: $P$ righe, ciascuna contenente 
    tre interi $S$, $E$, $C$ separati da uno spazio, dove $S$ è la 
    città di partenza, $E$ è la città di arrivo e $C$ è il costo 
    della tratta.
\end{itemize}

Le tratte sono \textbf{non orientate}: una tratta $(S, E, C)$ consente 
di viaggiare sia da $S$ a $E$ che da $E$ a $S$ con lo stesso costo $C$.

\subsection{Assunzioni}

\begin{itemize}
    \item $2 \leq N \leq 1000$
    \item $1 \leq P \leq 10000$
    \item $0 \leq S_i \leq N$ per ogni $i = 1, \ldots, P$
    \item $0 \leq E_i \leq N$ per ogni $i = 1, \ldots, P$
    \item $1 \leq C_i \leq 10000$ per ogni $i = 1, \ldots, P$
\end{itemize}

\subsection{Esempio di File di Input}

\begin{lstlisting}
4 6
0 1 4
0 2 1
1 3 2
1 4 12
2 3 4
3 4 5
\end{lstlisting}

\section{Output}

Il programma stampa su \texttt{stdout} una riga per ogni città 
raggiungibile dalla capitale (dalla città $1$ alla città $N$), 
nel seguente formato:

\begin{lstlisting}
S->...->E = totale-sconto=finale
\end{lstlisting}

dove:
\begin{itemize}
    \item \texttt{S->...->E} --- sequenza di città del percorso ottimo, 
    separate da \texttt{->};
    \item \texttt{totale} --- somma dei costi di tutte le tratte 
    del percorso;
    \item \texttt{sconto} --- costo della tratta più onerosa del percorso;
    \item \texttt{finale} --- costo effettivo dopo lo sconto 
    (\texttt{totale} $-$ \texttt{sconto}).
\end{itemize}

Se una città non è raggiungibile dalla capitale, viene stampato:

\begin{lstlisting}
Citta N: non raggiungibile
\end{lstlisting}

\subsection{Esempio di Output}

\begin{lstlisting}
0->1 = 4-4=0
0->2 = 1-1=0
0->2->3 = 5-4=1
0->1->4 = 16-12=4
\end{lstlisting}